\setcounter{section}{7}

  \zwischenueberschrift{Übungsaufgaben}

\inputaufgabe
{}
{

Zeige, dass zu zwei Punkten
\mathl{A,B \neq N,S}{} auf der Einheitssphäre $S^2$ die Differenz der in den beiden stereographischen Standardkarten genommenen Abständen, also
\mathkor {} {d( \alpha_1(A),\alpha_1(B))} {und} {d( \alpha_2(A),\alpha_2(B))} {,}
beliebig klein und beliebig groß sein kann.

}
{} {}
Die vorstehende Aufgabe zeigt, dass man über Karten im Allgemeinen keinen sinnvollen Abstandsbegriff auf einer Mannigfaltigkeit erhalten kann. Eine natürliche Metrik auf der Einheitssphäre ergibt sich durch die induzierte Metrik des $\R^3$ oder durch den geodätischen Abstand, siehe
[[Einheitssphäre/Metrik über Großkreis/Aufgabe|Aufgabe 7.16]].

\inputaufgabe
{}
{

Zeige, dass ein
\definitionsverweis {halboffenes Intervall}{}{}
\mathl{[a,b[}{} keine
\definitionsverweis {topologische Mannigfaltigkeit}{}{} ist.

}
{} {}

\inputaufgabe
{}
{

Es sei
\mathl{I=[0,1[}{} das
\zusatzklammer {nach oben} {} {}
halboffene Einheitsintervall und $S^1$ der
\definitionsverweis {Einheitskreis}{}{.}
Zeige, dass es eine
\definitionsverweis {bijektive}{}{}
\definitionsverweis {stetige Abbildung}{}{}
\maabbdisp {f} { I } { S^1
} {}
gibt, dass aber
\mathkor {} {I} {und} {S^1} {}
nicht
\definitionsverweis {homöomorph}{}{}
sind.

}
{} {}

\inputaufgabegibtloesung
{}
{

Zeige, dass die drei eindimensionalen Mannigfaltigkeiten
\mathlistdisp {S^1 = \left\{ (x,y) \in \R^2  \mid   \betrag { (x,y) } = 1         \right\}} {} {\R} {und} {\R \setminus \{0\}} {}
paarweise nicht homöomorph sind.

}
{} {}

\inputaufgabe
{}
{

Es sei $M$ der
\definitionsverweis {Graph}{}{}
der durch

\mavergleichskettedisp
{\vergleichskette
{ f(x)
}
{ =} {\begin{cases} \sin \frac{1}{x} \text{ für } x \neq 0\, , \\ 0 \text{ sonst}\, , \end{cases}
}
{ } {
}
{ } {
}
{ } {
}
}
{}{}{}
gegebenen Funktion
\maabb {f} {\R} {\R
} {.}
Zeige, dass $M$ keine
\definitionsverweis {topologische Mannigfaltigkeit}{}{}
ist.

}
{} {}

\inputaufgabe
{}
{

Zeige, dass ein
\definitionsverweis {offener Ball}{}{}

\mavergleichskette
{\vergleichskette
{ U \left(   P ,r  \right)
}
{ \subseteq }{ \R^m
}
{  }{
}
{    }{
}
{   }{
}
}
{}{}{}
$C^\infty$-\definitionsverweis {diffeomorph}{}{}
zum $\R^m$ ist.

}
{} {}

\inputaufgabegibtloesung
{}
{

Sei

\mathdisp {S= \left\{ P \in \R^3  \mid   \Vert { P } \Vert = 1         \right\}} {  }

die Einheitssphäre. Zu
\mathl{v=(a,b,c) \neq 0}{} ist

\mathdisp {E_v = \left\{ (x,y,z) \in \R^3  \mid  ax+by+cz = 0         \right\}} {  }

eine Ebene durch den Nullpunkt, die einen Großkreis
\zusatzklammer {einen \anfuehrung{Äquator}{}} {} {}
und zwei offene Halbsphären auf $S$ definiert.

a) Beschreibe zu
\mathl{v=(a,b,c) \neq 0}{} den zugehörigen Großkreis und die beiden Halbsphären mit Gleichungen bzw. mit Ungleichungen.

b) Zeige, dass man $S$ nicht mit drei offenen Halbsphären überdecken kann.

}
{} {}

\inputaufgabe
{}
{

Zeige, dass die offene Zylinderoberfläche
\mathl{S^1 \times {]0,1[}}{} zu
\mathl{S^1 \times \R}{,} zur punktierten Ebene
\mathl{\R^2 \setminus \{(0,0)\}}{} und zu
\mathl{S^2 \setminus \{N,S\}}{}
\definitionsverweis {homöomorph}{}{}
ist.

}
{} {}

\inputaufgabe
{}
{

Es sei
\mathl{]a,b[}{} ein
\definitionsverweis {offenes Intervall}{}{}
und
\maabbdisp {f} {]a,b[} {\R_{\geq 0}
} {}
eine
\definitionsverweis {stetige Funktion}{}{.}
Es sei $M$ die äußere Oberfläche des zugehörigen Rotationskörpers. Zeige, dass diese Menge bei

\mavergleichskette
{\vergleichskette
{ f
}
{ > }{ 0
}
{  }{
}
{    }{
}
{   }{
}
}
{}{}{}
eine zu einem offenen Zylindermantel homöomorphe Mannigfaltigkeit ist. Zeige ferner, dass keine Mannigfaltigkeit vorliegt, wenn $f$ sowohl Nullstellen als auch positive Werte besitzt.

}
{} {}

\inputaufgabe
{}
{

Zeige, dass eine
\definitionsverweis {topologische Mannigfaltigkeit}{}{}
genau dann
\definitionsverweis {zusammenhängend}{}{}
ist, wenn sie
\definitionsverweis {wegzusammenhängend}{}{}
ist.

}
{} {}

\inputaufgabe
{}
{

Es sei $M$ eine
\definitionsverweis {topologische Mannigfaltigkeit}{}{}
und es sei

\mathl{M= \bigcup_{i \in I} U_i}{} eine
\definitionsverweis {offene Überdeckung}{}{}
mit
\definitionsverweis {Karten}{}{}
\maabbdisp {\alpha_i} {U_i } { V_i
} {}
mit
\mathl{V_i \subseteq \R^n}{.} Zu
\mathl{i,j \in I}{} seien
\maabbdisp {\varphi_{ij} = \alpha_j \circ (\alpha_i)^{-1}} {V_i \cap \alpha_i(U_i \cap U_j)  } { V_j \cap \alpha_j(U_i \cap U_j)
} {}
die
\definitionsverweis {Übergangsabbildungen}{}{.}
Zeige, dass zu
\mathl{i,j,k \in I}{} die sogenannte \stichwort {Kozykelbedingung} {}
\zusatzklammer {auf welchen Teilmengen} {?} {}

\mavergleichskettedisp
{\vergleichskette
{ \varphi_{ij}
}
{ =} {\varphi_{kj} \circ \varphi_{ik}
}
{ } {
}
{ } {
}
{ } {
}
}
{}{}{}
gilt.

}
{} {}

\inputaufgabe
{}
{

Überprüfe die Kozykelbedingung
\zusatzklammer {siehe
[[Mannigfaltigkeit/Übergangsabbildungen/Kozykelbedingung/Aufgabe|Aufgabe 7.11]]} {} {}
für die Einheitssphäre $M=S^2$ und die drei stereographischen Projektionen vom Nordpol, vom Südpol und von
\mathl{P=(1,0,0)}{.}

}
{} {}

\inputaufgabegibtloesung
{}
{

Erstelle eine Animation, die die geometrischen Objekte aus
[[Einheitssphäre/Nordpol und (1,0,0)/Abstand längs Ebenenschnitt/Abhängigkeit vom Winkel/Aufgabe|Aufgabe 7.17]]
darstellt.

}
{} {}

\inputaufgabe
{}
{

Betrachte Geschenkpapier. Auf welche Arten kann man das Papier zerschneiden und/oder verkleben, so, dass eine zweidimensionale Mannigfaltigkeit entsteht. Sollte der Rand des Papiers dazu gehören oder nicht? Welche entstehenden Mannigfaltigkeiten sind zusammenhängend, welche kompakt? Wie entsteht ein Möbius-Band? Welche Möglichkeiten gibt es, wenn man endlich viele Ausnahmepunkte erlaubt, in denen keine Mannigfaltigkeitsstruktur vorliegt?

Wende die Theorie an, um möglichst originelle Verpackungen zu konstruieren. Verschnüre diese mit geeigneten eindimensionalen Mannigfaltigkeiten.

}
{} {}

  \zwischenueberschrift{Aufgaben zum Abgeben}

\inputaufgabe
{3}
{

Bestimme das
\definitionsverweis {Bild}{}{} der
\definitionsverweis {Großkreise}{}{}
durch die beiden Pole auf der
\definitionsverweis {Einheitssphäre}{}{} unter der
\definitionsverweis {stereographischen Projektion}{}{} vom Nordpol aus.

}
{} {}

\inputaufgabe
{5}
{

Zeige, dass auf der
\definitionsverweis {Einheitssphäre}{}{}

\mavergleichskette
{\vergleichskette
{ K
}
{ \subset }{ \R^3
}
{  }{
}
{    }{
}
{   }{
}
}
{}{}{}
durch folgende Zuordnung eine
\definitionsverweis {Metrik}{}{}
festgelegt wird. Für

\mavergleichskette
{\vergleichskette
{ P,Q
}
{ \in }{ K
}
{  }{
}
{    }{
}
{   }{
}
}
{}{}{}
ist
\mathl{d(P,Q)}{} die Länge des
\zusatzklammer {kürzeren} {} {} Verbindungsweges von $P$ nach $Q$ auf dem durch diese Punkte festgelegten
\definitionsverweis {Großkreis}{}{}
\zusatzklammer {berücksichtige auch die Fälle

\mavergleichskette
{\vergleichskette
{ P
}
{ = }{ Q
}
{  }{
}
{    }{
}
{   }{
}
}
{}{}{}
und $P,Q$
\definitionsverweis {antipodal}{}{}} {} {.}

}
{} {}

\inputaufgabe
{8}
{

Wir fixieren die beiden Punkte $N=(0,0,1)$ und $P=(1,0,0)$ auf der
\definitionsverweis {Einheitssphäre}{}{} $K$.  Es sei $G$ die Verbindungsgerade und es sei $H$ die zu $G$ senkrechte Ebene durch $N$. Führe auf $H$ einen parametrisierten Einheitskreis $E$ mit $N$ als Mittelpunkt ein. Bestimme zu $S \in E$ die Länge des
\zusatzklammer {kürzeren} {} {} Weges von $N$ nach $P$ auf demjenigen Kreis, der durch den Schnitt von $K$ mit der durch
\mathkor {} {N,P} {und} {S} {} gegebenen Ebene festgelegt ist.

}
{} {}

\inputaufgabe
{6}
{

Man gebe eine
\definitionsverweis {injektive}{}{}
\definitionsverweis {stetige Abbildung}{}{}
\maabbdisp {\varphi} {\R} {S^2
} {,}
die
\zusatzklammer {als Abbildung nach $\R^3$} {} {}
\definitionsverweis {rektifizierbar}{}{}
ist und unendliche
\definitionsverweis {Länge}{}{}
besitzt, und für die
\mathkor {} {\operatorname{lim}_{   t  \rightarrow  \infty  } \, \varphi (t) = N} {und} {\operatorname{lim}_{   t  \rightarrow  - \infty  } \, \varphi (t) = S} {}
gilt.

}
{} {}

\inputaufgabe
{4}
{

Zeige, dass das Achsenkreuz keine
\definitionsverweis {topologische Mannigfaltigkeit}{}{}
ist.

}
{} {}

 [[Kategorie:Latexseite]]
